% from ega0-15-fr.tex
\oldpage[0\textsubscript{IV-1}]{21}
de la façon suivante~: pour tout ouvert $U$ de $X$ (ou d'une base de la
topologie de $X$), on considère l'homomorphisme (15.1.1.1)
\[
  \varphi_{k,U}:\gr_0(\Gamma(U,\mathscr{F}))\otimes_{\mathbb Z}S_k
  \longrightarrow\gr_k(\Gamma(U,\mathscr{F}))
\]
défini par les $n$ sections $f_i|U\in\Gamma(U,\mathcal{O}_X)$, le
$\Gamma(U,\mathcal{O}_X)$-module $\Gamma(U,\mathscr{F})$ étant filtré par les
sous-modules $(\Gamma(U,\mathscr{J}))^k\Gamma(U,\mathscr{F})$. Le
$\mathcal{O}_X$-Module $\gr_k(\mathscr{F})$ étant associé au préfaisceau
$U\mapsto\gr_k(\Gamma(U,\mathscr{F}))$ pour tout $k$, les $\varphi_{k,U}$
définissent un homomorphisme de $\mathcal{O}_X$-Modules (15.2.1.1). En vertu
des propriétés d'exactitude des limites inductives de modules et de leur
commutation aux produits tensoriels, il résulte de ce qui précède que pour
tout $x\in X$, l'homomorphisme $(\varphi_k)_x$ des fibres des deux membres de
(15.2.1.1) au point $x$ s'identifie à l'homomorphisme de (15.1.1)
\[
\label{0.15.2.1.2-fr}
  \gr_0(\mathscr{F}_x)\otimes_{\mathbb Z}S_k
  \longrightarrow\gr_k(\mathscr{F}_x)
\tag{15.2.1.2}
\]
défini par la suite $((f_i)_x)_{1\leq i\leq n}$ d'éléments de
$\mathcal{O}_x$.
\end{env}

\begin{env}[15.2.2]
\label{0.15.2.2-fr}
Les notations étant celles de (15.2.1), nous dirons que la suite
$(f_i)_{1\leq i\leq n}$ est $\mathscr{F}$-régulière si, en désignant par
$f_j\mathscr{F}$ le sous-$\mathcal{O}_X$-Module de $\mathscr{F}$ image de
l'endomorphisme de $\mathscr{F}$ défini par $f_j$, l'endomorphisme de
$\mathscr{F}/(\sum_{j=1}^{i-1}f_j\mathscr{F})$ défini par $f_i$ est injectif
pour $1\leq i\leq n$~; nous dirons que la suite $(f_i)$ est
$\mathscr{F}$-quasi-régulière si les homomorphismes (15.2.1.1) sont injectifs.
Il résulte de ces définitions et des remarques de (15.2.1) que pour que la
suite $(f_i)$ soit $\mathscr{F}$-régulière (resp.
$\mathscr{F}$-quasi-régulière), il faut et il suffit que, pour tout $x\in X$,
la suite $((f_i)_x)_{1\leq i\leq n}$ soit $\mathscr{F}_x$-régulière (resp.
$\mathscr{F}_x$-quasi-régulière). Toute suite $\mathscr{F}$-régulière est donc
$\mathscr{F}$-quasi-régulière (15.1.9)~; la réciproque est vraie si
$\mathscr{F}$ est un $\mathcal{O}_X$-Module de type fini, si les
$\mathcal{O}_x$ sont des anneaux noethériens pour tout $x\in X$ et si en outre
les $(f_i)_x$ appartiennent au radical de $\mathcal{O}_x$ pour tout $x\in X$
(15.1.11).

Si $X=\operatorname{Spec}(A)$ est un schéma affine, et
$\mathscr{F}=\widetilde{M}$, il
résulte de (I, 1.3) que pour que la suite $(f_i)$ soit
$\mathscr{F}$-régulière (resp. $\mathscr{F}$-quasi-régulière), il faut et il
suffit qu'elle soit $M$-régulière (resp. $M$-quasi-régulière).

Lorsque $\mathscr{F}=\mathcal{O}_X$, on supprimera la mention de
$\mathscr{F}$ dans les définitions précédentes.
\end{env}

\begin{remark}[15.2.3]
\label{0.15.2.3-fr}
Lorsque $X=\operatorname{Spec}(A)$ et
$\mathfrak{J}=\sum_{i=1}^n f_iA$, la suite $(f_j)$ est
quasi-régulière pour tout $A$-module $M$ dans tout ouvert $U$ ne rencontrant
pas $V(\mathfrak{J})$, car en tout point $x\in U$, on a
$\mathfrak{J}_x=\mathcal{O}_x$ et les deux membres de (15.2.1.2) sont nuls. La
notion de suite $\widetilde{M}$-régulière n'a donc d'intérêt qu'au voisinage
des points de $V(\mathfrak{J})$.
\end{remark}

\begin{proposition}[15.2.4]
\label{0.15.2.4-fr}
Soient $X$ un espace annelé, $\mathscr{F}$ un $\mathcal{O}_X$-Module cohérent,
$x$ un point de $X$, $(f_i)_{1\leq i\leq n}$ une suite d'éléments de
$\Gamma(X,\mathcal{O}_X)$. Si la suite $((f_i)_x)_{1\leq i\leq n}$ est
$\mathscr{F}_x$-régulière, il existe un voisinage ouvert $U$ de $x$ tel que la
suite $(f_i|U)_{1\leq i\leq n}$ soit $(\mathscr{F}|U)$-régulière.
\end{proposition}

\begin{proof}
En effet, pour $1\leq i\leq n$, le $\mathcal{O}_X$-Module
$\mathscr{F}/(\sum_{j=1}^{i-1}f_j\mathscr{F})$ est cohérent
($0_{\mathrm I}$, 5.3.4), donc il en est de même du noyau de
l'endomorphisme de ce $\mathcal{O}_X$-Module défini par $f_i$
($0_{\mathrm I}$, 5.3.4)~; comme la fibre au point $x$ de ce noyau est nulle,
il en est de même de sa restriction à un voisinage de $x$
($0_{\mathrm I}$, 5.2.2).
\end{proof}

\begin{proposition}[15.2.5]
\label{0.15.2.5-fr}
